% Part: first-order-logic
% Chapter: completeness
% Section: henkin-expansion

\documentclass[../../../include/open-logic-section]{subfiles}

\begin{document}

\olfileid{fol}{com}{hen}

\olsection{Ekspansi Henkin}

\begin{explain}
Salah siji tantangan nalika mbuktekake teorema kelengkapan yaiku model
kang dibangun saka himpunan konsisten jangkep~$\Gamma$ kudu ndadekake
kabeh !!{formula}s kang dikuantifikasi ing~$\Gamma$ bener. Kanggo
njamin iki, kita nggunakake siasat saka Leon Henkin. Intine, siasat
iki nggedhekake basa nganggo tanpa winates akeh !!{constant}s lan
nambah, kanggo saben !!{formula} kanthi siji !!{variable} bebas
$!A(x)$, formula awujud
\iftag{prvEx}
      {$\lexists[x][!A(x)] \lif !A(c)$}
      {$\lnot\lforall[x][!A(x)] \lif \lnot !A(c)$},
kanthi $c$ minangka salah siji !!{constant}s anyar. Nalika kita
mbangun !!{structure} kang nyukupi~$\Gamma$, iki bakal njamin yen saben
\iftag{prvEx}
{ukara eksistensial kang bener nduweni saksi}
{ukara universal kang luput nduweni conto panolak}
ing antarane konstanta anyar kasebut.
\end{explain}

\begin{prop}
\ollabel{prop:lang-exp}
Yen $\Gamma$ konsisten ing $\Lang L$ lan $\Lang L'$ dipikolehi saka
$\Lang L$ kanthi nambah himpunan kang !!a{denumerable} saka !!{constant}s anyar
$\Obj d_0$, $\Obj d_1$, \dots, mula $\Gamma$ konsisten ing~$\Lang L'$.
\end{prop}

\begin{defn}[Himpunan jenuh]
Himpunan $\Gamma$ saka !!{formula}s ing basa $\Lang {L}$ diarani
\emph{jenuh} yen lan mung yen kanggo saben !!{formula}~$!A(x) \in
\Frm[L]$ kanthi siji !!{variable} bebas~$x$, ana !!a{constant}~$c \in
\Lang{L}$ saengga
\iftag{prvEx}
      {$\lexists[x][!A(x)] \lif !A(c) \in \Gamma$}
      {$\lnot\lforall[x][!A(x)] \lif \lnot !A(c) \in \Gamma$}.
\end{defn}

Definisi ing ngisor iki bakal digunakake ing bukti teorema sabanjure.

\begin{defn}
\ollabel{defn:henkin-exp}
Upamana $\Lang L'$ kaya ing \olref{prop:lang-exp}. Tetepna enumerasi
$!A_0(x_0)$, $!A_1(x_1)$, \dots saka kabeh !!{formula}s~$!A_i(x_i)$
ing~$\Lang L'$ kang ngemot siji variabel ($x_i$) kanthi bebas. Kita
nemtokake !!{sentence}s~$!D_n$ kanthi induksi marang~$n$.

Upamana $c_0$ minangka !!{constant} kapisan ing antarane $\Obj d_i$
kang ditambahake marang~$\Lang{L}$ lan ora muncul ing~$!A_0(x_0)$.
Kanthi nganggep $!D_0$, \dots,~$!D_{n-1}$ wis ditemtokake, upamana
$c_n$ minangka sing kapisan ing antarane !!{constant}s anyar~$\Obj
d_i$ kang ora muncul ing $!D_0$, \dots,~$!D_{n-1}$ lan uga ora muncul
ing~$!A_n(x_n)$.

Saiki upamana $!D_{n}$ minangka !!{formula} \iftag{prvEx}
{$\lexists[x_{n}][!A_{n}(x_{n})] \lif
  !A_{n}(c_{n})$}{$\lnot\lforall[x_{n}][!A_{n}(x_{n})] \lif \lnot
  !A_{n}(c_{n})$}.
\end{defn}

\begin{lem}
\ollabel{lem:henkin}
Saben himpunan konsisten~$\Gamma$ bisa diambakake dadi himpunan
konsisten jenuh~$\Gamma'$.
\end{lem}

\begin{proof}
Saka himpunan ukara konsisten~$\Gamma$ ing basa~$\Lang{L}$,
gedhekna basane kanthi nambah himpunan kang !!a{denumerable} saka
!!{constant}s anyar kanggo mbentuk~$\Lang{L'}$. Miturut
\olref{prop:lang-exp}, $\Gamma$ isih konsisten ing basa kang luwih
sugih iki. Sabanjure, upamana $!D_i$ kaya ing
\olref{defn:henkin-exp}. Tetepna
\begin{align*}
\Gamma_0 & = \Gamma \\
\Gamma_{n+1} & = \Gamma_n \cup \{!D_n \}
\end{align*}
yaiku, $\Gamma_{n+1} = \Gamma \cup \{ !D_0, \dots, !D_n \}$, lan
tetepna $\Gamma' = \bigcup_{n} \Gamma_n$. Cetha yen $\Gamma'$ jenuh.

Yen $\Gamma'$ ora konsisten, kanggo sawijine $n$, $\Gamma_n$ mesthi
ora konsisten (Gladhen: terangna sebabe). Mula, kanggo nuduhake yen
$\Gamma'$ konsisten, cukup dituduhake kanthi induksi marang~$n$ yen
saben himpunan~$\Gamma_n$ konsisten.

Basis induksine mung pratelan yen $\Gamma_0 = \Gamma$ konsisten, kang
minangka hipotesis teorema. Kanggo langkah induksi, upamana
$\Gamma_{n}$ konsisten nanging $\Gamma_{n+1} = \Gamma_n \cup
\{!D_n\}$ ora konsisten. Elinga yen $!D_n$~yaiku
\iftag{prvEx}
{$\lexists[x_{n}][!A_{n}(x_n)] \lif !A_{n}(c_{n})$}
{$\lnot\lforall[x_{n}][!A_{n}(x_n)] \lif \lnot !A_{n}(c_{n})$},
kanthi $!A_n(x_n)$ minangka !!a{formula} saka $\Lang{L'}$ kang mung
nduweni variabel~$x_n$ kanthi bebas. Miturut cara kita milih~$c_n$
(delengen \olref{defn:henkin-exp}), $c_n$ ora muncul ing~$!A_n(x_n)$
utawa ing~$\Gamma_n$.

Yen $\Gamma_n \cup \{!D_n\}$ ora konsisten, mula $\Gamma_n
\Proves \lnot !D_n$, lan mulane rong pratelan ing ngisor iki bener:
\iftag{prvEx}{
\[
\Gamma_n \Proves \lexists[x_n][!A_n(x_n)]
\qquad
\Gamma_n \Proves \lnot !A_n(c_n)
\]}{
\[
\Gamma_n \Proves \lnot\lforall[x_n][!A_n(x_n)]
\qquad
\Gamma_n \Proves !A_n(c_n)
\]}
Amarga $c_n$ ora muncul ing
$\Gamma_n$ utawa ing~$!A_n(x_n)$,
\tagrefs{prfAX/{fol:axd:qpr:thm:strong-generalization},prfSC/{fol:seq:qpr:thm:strong-generalization},prfND/{fol:ntd:qpr:thm:strong-generalization},prfTab/{fol:tab:qpr:thm:strong-generalization}} bisa ditrapake.
Saka \iftag{prvEx}
{$\Gamma_n \Proves \lnot !A_n(c_n)$}
{$\Gamma_n \Proves !A_n(c_n)$},
kita entuk
\iftag{prvEx}
{$\Gamma_n \Proves \lforall[x_n][\lnot !A_n(x_n)]$}
{$\Gamma_n \Proves \lforall[x_n][!A_n(x_n)]$}.
Mula kita nduweni kalorone
\iftag{prvEx}
{$\Gamma_n \Proves \lexists[x_n][!A_n(x_n)]$ lan
$\Gamma_n \Proves \lforall[x_n][\lnot !A_n(x_n)]$}
{$\Gamma_n \Proves \lnot\lforall[x_n][!A_n(x_n)]$ lan
$\Gamma_n \Proves \lforall[x_n][!A_n(x_n)]$},
dadi $\Gamma_n$ dhewe ora konsisten.
\iftag{prvEx}
{(Cathetan:
\iftag{prvAll}
{$\lforall[x_n][\lnot !A_n(x_n)] \Proves
  \lnot\lexists[x_n][!A_n(x_n)]$}
{$\lforall[x_n][\lnot !A_n]$ ditetepake minangka
  $\lnot\lexists[x_n][\lnot\lnot !A_n(x_n)]$ lan
  $\lnot\lexists[x_n][\lnot\lnot !A_n(x_n)] \Proves
  \lnot\lexists[x_n][!A_n(x_n)]$}.)}{}
Iki cengkah: $\Gamma_n$ dikira konsisten. Mula $\Gamma_n \cup \{
!D_n\}$ konsisten.
\end{proof}

\begin{explain}
Saiki bakal dituduhake yen himpunan konsisten kang \emph{jangkep} lan
jenuh nduweni sipat yen \iftag{prvAll}{himpunan iku ngemot
  !!{sentence} kang dikuantifikasi universal yen lan mung yen ngemot
  kabeh conto ukarane}{}\iftag{defAll,defEx}{}{ lan
  }\iftag{prvEx}{himpunan iku ngemot !!{sentence} kang dikuantifikasi
  eksistensial yen lan mung yen ngemot paling ora siji conto ukarane}{}.
Iki bakal digunakake kanggo nuduhake yen !!{structure} kang bakal
diasilake saka himpunan kang jangkep, konsisten, lan jenuh ndadekake
kabeh ukara kang dikuantifikasi ing himpunan kasebut bener. Cathetan koreksi sumber OLPL-740: cabang andharan eksistensial nggunakake tag prvEx, cocog karo proposisi lan bukti jejer; sumber salah nulis prvAll ing cabang iki.
\end{explain}

\begin{prop}\ollabel{prop:saturated-instances}
Upamana $\Gamma$ iku jangkep, konsisten, lan jenuh.
\begin{tagenumerate}{prvEx,prvAll}
\tagitem{prvEx}{$\lexists[x][!A(x)] \in \Gamma$ yen lan mung yen
  $!A(t) \in \Gamma$ kanggo paling ora siji term katutup~$t$.}{}
\tagitem{prvAll}{$\lforall[x][!A(x)] \in \Gamma$ yen lan mung yen
  $!A(t) \in \Gamma$ kanggo kabeh term katutup~$t$.}{}
\end{tagenumerate}
\end{prop}

\begin{proof}
  \begin{tagenumerate}{prvEx,prvAll}
  \tagitem{prvEx}{%
    \iftag{probEx}{Gladhen.}{Kaping pisan, upamana
      $\lexists[x][!A(x)] \in \Gamma$. Amarga $\Gamma$ jenuh,
      $(\lexists[x][!A(x)] \lif !A(c)) \in \Gamma$ kanggo sawijine
      !!{constant}~$c$. Miturut
      \tagrefs{prfAX/{fol:axd:ppr:prop:provability-lif},prfSC/{fol:seq:ppr:prop:provability-lif},prfND/{fol:ntd:ppr:prop:provability-lif},prfTab/{fol:tab:ppr:prop:provability-lif}}, butir~(1),
      lan \olref[ccs]{prop:ccs}\olref[ccs]{prop:ccs-prov-in}, $!A(c)
      \in \Gamma$.

    Kanggo arah liyane, sipat jenuh ora dibutuhake: upamana
    $!A(t) \in \Gamma$. Banjur $\Gamma \Proves
    \lexists[x][!A(x)]$ miturut
      \tagrefs{prfAX/{fol:axd:qpr:prop:provability-quantifiers},prfSC/{fol:seq:qpr:prop:provability-quantifiers},prfND/{fol:ntd:qpr:prop:provability-quantifiers},prfTab/{fol:tab:qpr:prop:provability-quantifiers}}, butir~(1). Miturut
    \olref[ccs]{prop:ccs}\olref[ccs]{prop:ccs-prov-in},
    $\lexists[x][!A(x)] \in \Gamma$.}}{}
  \tagitem{prvAll}{%
    \iftag{probAll}{Gladhen.}{Upamana $!A(t) \in \Gamma$ kanggo
      kabeh term katutup~$t$. Kanggo nggayuh cengkah, upamana
      $\lforall[x][!A(x)] \notin \Gamma$. Amarga $\Gamma$ jangkep,
      $\lnot\lforall[x][!A(x)] \in \Gamma$. Amarga jenuh,
      \iftag{prvEx}{$(\lexists[x][\lnot !A(x)] \lif \lnot !A(c)) \in
        \Gamma$}{$(\lnot\lforall[x][!A(x)] \lif \lnot !A(c)) \in
        \Gamma$} kanggo sawijine !!{constant}~$c$. Miturut hipotesis,
      amarga $c$ minangka term katutup, $!A(c) \in \Gamma$. Nanging
      iki bakal ndadekake $\Gamma$ ora konsisten. (Gladhen: wenehana
      !!{derivation} kang nuduhake yen
      \[
      \lnot \lforall[x][!A(x)],
      \iftag{prvEx}{\lexists[x][\lnot !A(x)] \lif \lnot
        !A(c)}{\lnot\lforall[x][!A(x)] \lif \lnot
        !A(c)}, !A(c)
      \]
      ora konsisten.)

      Kanggo arah kosok baline, kita ora mbutuhake sipat jenuh: upamana
      $\lforall[x][!A(x)] \in \Gamma$. Banjur $\Gamma \Proves !A(t)$
      miturut
      \tagrefs{prfAX/{fol:axd:qpr:prop:provability-quantifiers},prfSC/{fol:seq:qpr:prop:provability-quantifiers},prfND/{fol:ntd:qpr:prop:provability-quantifiers},prfTab/{fol:tab:qpr:prop:provability-quantifiers}},
      butir~(2). Kita entuk $!A(t) \in \Gamma$ miturut
      \olref[ccs]{prop:ccs}.}}{}
  \end{tagenumerate}
\end{proof}

\end{document}
